

\newpage
\chapter{Introduction}
\mbox{}\vspace{-\baselineskip}

Color transparency (CT) is an interesting prediction of Quantum Chromodynamics (QCD) that describes a situation occurring in hard interactions with nuclei at high energies. Specifically, color transparency refers to the situation when hadrons propagate through nuclear medium with less interaction than was expected. This phenomenon is attributed to the fact that at high momentum transfers hadrons are thought to form so-called point-like configurations (PLCs) that are color-neutral ($q\bar q$ or $qqq$). For such configurations, cancellation of the QCD color fields takes place, and as a result, PLCs interact with the nuclear medium with the reduced cross section. Therefore, color transparency is really about how a medium can be transparent to objects without color~\cite{Holly_color_t,Dutta_color_t}.



Over the past decades, numerous experimental efforts have aimed to observe and measure CT effects in electron and hadron scattering off different targets. For the $q\bar q$ configuration, clear indications of the enhanced nuclear transparency were observed in experiments of meson electroproduction at JLab (see Ref.~\cite{Clasie:2007aa} for $\pi^+$-production off ${}^{2}\text{H}$, ${}^{12}\text{C}$, ${}^{27}\text{Al}$, ${}^{63}\text{Cu}$, and ${}^{197}\text{Au}$ at 1.1~GeV$^{2}<Q^{2}<$~4.7~GeV$^{2}$ and Ref.~\cite{ElFassi:2012hsy} for $\rho^{0}$-production off ${}^{12}\text{C}$ and ${}^{56}\text{Fe}$ at 1.0~GeV$^{2}<Q^{2}<$~2.2~GeV$^{2}$). For the $qqq$ configuration, CT effects have also been searched for in reactions of hard $p$ knock-out off different targets (${}^{2}\text{H}$, ${}^{12}\text{C}$, ${}^{63}\text{Cu}$) in quasielastic electron scattering at JLab~\cite{Garrow:2001di,HallC:2020ijh} for $Q^{2}$ up to 14~GeV$^{2}$. However, in these experiments, no evidence of color transparency were revealed. A possible explanation for that could be that for a proton, a PLC system may not form, or may expand too quickly in the investigated $Q^2$ range, meaning that larger $Q^2$ values might be needed to form a PLC~\cite{Brodsky:2022bum,Frankfurt:2022cyk}.



Meanwhile, some indications of color transparency were observer for the $qqq$ configuration in the $(p,2p)$ quasi-elastic scattering off carbon with the EVA detector at the Brookhaven AGS at near 90$^{o}$ $pp$ center of mass~\cite{PhysRevLett.87.212301}. As follows from Ref.~\cite{PhysRevLett.87.212301}, for such a process, one can expect color transparency effects to become discernible starting from $Q^{2}\gtrsim$~5~GeV$^{2}$. 

In summary, evidence for color transparency for protons has been much less conclusive compared to mesons, and for these reason, more measurements for the $p$ knock-out process are needed for higher values of $Q^2$, to better understand CT effects for the $qqq$ configuration. Such measurements can be performed at the Nuclotron-based Ion Collider fAcility (NICA) at JINR in Dubna.



In general, NICA beam time is expected to be shared between the two major experiments, each employing a corresponding detector facility located at one of the interaction points at the NICA collider rings: MultiPurpose Detector (MPD) and Spin Physics Detector (SPD). The two experiments are expected to operate in succession, and for each of them, the operation time will be subdivided into different ``stages", where each stage will have different beam requirements, run conditions, and detector setups, and thus essentially different physics programs\cite{TDR_SPD,SPDproto:2021hnm,SPD_1st_stage,Arbuzov:2020cqg}. 

Measurements of CT effects are proposed to be performed with the SPD detector during the first stage of the NICA collider operation, for which NICA is expected to deliver high-intensity polarized proton ($p$--$p$) and deuteron ($d$--$d$) beams at a collision energy up to 27~GeV and a luminosity up to 10$^{32}$ cm$^{-2}$ s$^{-1}$\cite{TDR_SPD,SPDproto:2021hnm,SPD_1st_stage}.

Due to the fact that asymmetric beam mode ($p$--$d$) is not expected to be available during the first stage of the SPD experiment, it was suggested in Ref.~\cite{Larionov:2025equ} to use $d$--$d$ interactions to study the process of $p$ knock-out from a deuteron in hard $p$--$d$ collisions, assuming the neutron in one of the beam-deuterons to be a spectator to the investigated process.

This report presents a thorough investigation of the prospects of measuring color transparency effects in the upcoming NICA SPD experiment. Chapter~\ref{Chapter_observables} gives a brief and general overview of the observables sensitive to color transparency effects. Chapter~\ref{Chap:kin} describes the kinematics of the $pd\rightarrow ppn$ process and the related kinematic variables. 

To examine the prospects of the color transparency measurements in the NICA SPD experiment, the Monte Carlo simulation was performed using the event generator for the process of the proton knock-out from a deuteron in $d$--$d$ collisions. The event generator is based on the theoretical cross section approximation obtained in the framework of the generalized eikonal approach supplemented by the quantum diffusion model of CT~\cite{Frankfurt:1996uz,PhysRevC.107.014605,Larionov:2025equ,Farrar:1988me}. Chapter~\ref{chapter_rate} presents important details on the theoretical predictions for the $pd\rightarrow ppn$ cross sections and provides an estimation of the event rates of the proposed experiment based on the generated Monte Carlo events of the investigated process. The performed event rate estimation allowed us to develop the favorable kinematic settings for the proposed measurements, which are also described in Chapter~\ref{chapter_rate}. Chapter~\ref{chapter_generated} presents examination of the angular and momentum distributions of the generated Monte Carlo events.


Chapter~\ref{chap:event_selection} briefly outlines the overall setup of the SPD experiment and also the starting SPD configuration for the first stage of the NICA collider operation~\cite{TDR_SPD}. This chapter also describes the specificities of the proton and neutron identification during the first stage, presents the proposed event selection procedure for the investigated $pd\rightarrow ppn$ process, and provides some insights into potential challenges that may arise during the measurement itself and on the subsequent data analysis stage.


